% from ega0-16-fr.tex
\oldpage[0\textsubscript{IV-1}]{28}

\begin{proof}
En effet, $M/\mathfrak{r}^nM$ et
$\widehat{M}/\mathfrak{r}^n\widehat{M}$ sont isomorphes, donc
$d(\widehat{M})=d(M)$.
\end{proof}

\begin{corollary}[16.2.5]
\label{0.16.2.5-fr}
La dimension d'un anneau semi-local noethérien $A$ est finie et égale au
nombre minimum d'éléments du radical de $A$ engendrant un idéal de définition.
\end{corollary}

\begin{proof}
C'est l'égalité $s(M)=\dim(M)$ pour $M=A$.
\end{proof}

En particulier~:

\begin{corollary}[16.2.6]
\label{0.16.2.6-fr}
Soient $A$ un anneau local noethérien, $\mathfrak{m}$ son idéal maximal,
$k=A/\mathfrak{m}$ son corps résiduel~; on a
\begin{equation}
\label{0.16.2.6.1-fr}
\tag{16.2.6.1}
  \dim(A)\leq\operatorname{rg}_k(\mathfrak{m}/\mathfrak{m}^2).
\end{equation}
\end{corollary}

\begin{proof}
En effet, on sait que si les $x_i$ $(1\leq i\leq n)$ sont des éléments de
$\mathfrak{m}$ dont les classes mod.~$\mathfrak{m}^2$ forment une base du
$k$-espace vectoriel $\mathfrak{m}/\mathfrak{m}^2$, les $x_i$ engendrent
$\mathfrak{m}$ (Bourbaki, \emph{Alg.~comm.}, chap.~II, \S~3, no~2, cor.~2 de
la prop.~4).
\end{proof}

\begin{proposition}[16.2.7]
\label{0.16.2.7-fr}
Soient $k$ un corps, $B$ une $k$-algèbre graduée de type fini à degrés
positifs engendrée par ses éléments homogènes de degré 1 et telle que $B_0=k$~;
soit $E$ un $B$-module gradué de type fini, de sorte (III, 2.5.4) que, pour $n$
assez grand, $\operatorname{rg}_k(E_n)=r(n)$ est de la forme
$\Phi(n)-\Phi(n-1)$, où $\Phi$ est un polynôme en $n$ de degré $d$. Alors il
existe dans $B$ une famille $(t_i)_{1\leq i\leq d}$ de $d$ éléments homogènes
de degrés $\geq1$ tels que $E/(\sum_i t_iE)$ soit de rang fini sur $k$.
\end{proposition}

\begin{proof}
Soit $\mathfrak{q}$ l'idéal maximal $\bigoplus_{n\geq1}B_n$ de $B$, et, pour
tout $n$, soit $E'_n$ le sous-$B$-module $\bigoplus_{h\geq n+1}E_h$ de $E$.
Tout élément de $B-\mathfrak{q}$ détermine une homothétie injective dans le
$B$-module artinien $E/E'_n$, et cette homothétie est donc bijective
(Bourbaki, \emph{Alg.}, chap.~VIII, \S~2, no~2, lemme~3), ce qui entraîne que
$(E/E'_n)_{\mathfrak{q}}$ s'identifie à $E/E'_n$~; comme le corps résiduel de
$B_{\mathfrak{q}}$ est $k$, on déduit de l'hypothèse que
$\operatorname{long}_{B_{\mathfrak{q}}}((E/E'_n)_{\mathfrak{q}})$ est un
polynôme en $n$ de degré $d$. D'autre part, comme $E$ est un $B$-module de
type fini et que $B$ est engendré par ses éléments homogènes de degré 1, la
filtration $(E'_n)$ sur $E$ est $\mathfrak{q}$-bonne (Bourbaki,
\emph{Alg.~comm.}, chap.~III, \S~1, no~3, lemme~1), donc la filtration
$((E'_n)_{\mathfrak{q}})$ sur $E_{\mathfrak{q}}$ est
$\mathfrak{q}B_{\mathfrak{q}}$-bonne. Puisque $B_{\mathfrak{q}}$ est un anneau
local noethérien, on conclut de (16.2.3) et de (16.2.2.1) que l'on a
$\dim(E_{\mathfrak{q}})=d$. Raisonnons alors par récurrence sur $d$, le lemme
étant évident si $d=0$. Supposons $d\geq1$, et observons que les idéaux premiers
minimaux $\mathfrak{p}_i$ $(1\leq i\leq r)$ dans
$\operatorname{Ass}(E)$ sont \emph{gradués} (Bourbaki, \emph{Alg.~comm.},
chap.~IV, \S~3, no~1, prop.~1), et que $\mathfrak{q}$ est distinct de la
réunion des $\mathfrak{p}_i$ puisque $d\geq1$~; il y a donc un élément
\emph{homogène} $t_1\in\mathfrak{q}$ n'appartenant à aucun des
$\mathfrak{p}_i$ (Bourbaki, \emph{Alg.~comm.}, chap.~III, \S~1, no~4,
prop.~8). Comme les idéaux premiers $(\mathfrak{p}_i)_{\mathfrak{q}}$ sont les
éléments minimaux de $\operatorname{Ass}(E_{\mathfrak{q}})$ (Bourbaki,
\emph{Alg.~comm.}, chap.~IV, \S~1, no~2, prop.~5), on a
$\dim(E_{\mathfrak{q}}/t_1E_{\mathfrak{q}})=d-1$ (16.3.4). Il suffit alors
d'appliquer l'hypothèse de récurrence au $B$-module gradué $E/t_1E$ pour
obtenir la conclusion.
\end{proof}

\subsection{Systèmes de paramètres dans un anneau local noethérien}
\label{subsection:0.16.3-fr}

\begin{proposition}[16.3.1]
\label{0.16.3.1-fr}
Soient $A$ un anneau noethérien, $\mathfrak{p}$ un idéal premier de $A$, $n$ un
entier. Les conditions suivantes sont équivalentes~:
\begin{enumerate}
  \item[$a)$] $\operatorname{ht}(\mathfrak{p})=
  \dim(A_{\mathfrak{p}})\leq n$~;
  \item[$b)$] Il existe $n$ éléments $x_i$ de $A$ $(1\leq i\leq n)$ tels que
  $\mathfrak{p}$ soit minimal dans l'ensemble des idéaux premiers contenant
  l'idéal $\mathfrak{a}$ engendré par les $x_i$ (autrement dit
  $V(\mathfrak{p})$ est une composante irréductible de $V(\mathfrak{a})$).
\end{enumerate}
\end{proposition}

\begin{proof}
La condition $b)$ entraîne que $\mathfrak{a}A_{\mathfrak{p}}$ est un idéal de
définition de $A_{\mathfrak{p}}$, d'où $a)$ en vertu de (16.2.5).
Inversement, si $a)$ est vérifiée, il existe un idéal de définition
$\mathfrak{b}$ de $A_{\mathfrak{p}}$, engendré par $n$ éléments $x_i/s$, avec
$s\in A-\mathfrak{p}$. En vertu de la correspondance biunivoque entre idéaux
premiers de $A_{\mathfrak{p}}$ et idéaux premiers de $A$ contenus dans
$\mathfrak{p}$, l'idéal $\mathfrak{a}$ de $A$ engendré par les $x_i$ vérifie
$b)$.
\end{proof}

